\section{Quantitative estimates for the observations}\label{sec:observations}

The coarse cap estimates of Proposition~\ref{prop:edge-cap-inputs} locate
the relevant fields and bound their covariance by a constant times $p^{-2}$.
For the time normalization in Theorem~\ref{thm:main}, we shall need a relative
covariance error that is integrable in logarithmic scale. The purpose of
this section is to obtain the averaged estimates from which that improvement
will follow. There are two different sources of smallness: a two-replica
comparison at the smallest observation scale, and an integrated likelihood
estimate at the larger scales.

We retain the notation and conditional-on-spectrum probability
$\mathbf P$ of Lemma~\ref{lem:edge-biased-law}. Thus $\mathbf P$ includes the
partition-function bias on the eigenframe and the ordinary Gibbs observation
path given that frame; $Q_\mu$ denotes the latter path law for fixed disorder.
Expectations with respect to these laws are denoted by $\mathbf E$ and
$Q_\mu[\,\cdot\,]$. Constants in the present section are uniform on each of
the tight spectral sets described in Remark~\ref{rem:edge-tight-convergence}.
They may depend on that set, on a fixed moment order, and on fixed field or
scale ratios. A bound $C\exp(-n^\xi)$ always has some $\xi>0$; decreasing
$\xi$ absorbs polynomial factors and the $O(\log n)$ scale intervals used
below. Recall that
\[
 r=n^{-1/3+\delta},\qquad \delta=10^{-4},\qquad k_p=np^3,
 \qquad M_p=(np)^{1/2}.
\]
All estimates involving a fixed multiplicative enlargement of the interval
of scales have the same meaning, provided its upper endpoint is sufficiently
small. In particular, $k_p$ grows as a positive power of $n$ throughout the
intervals considered here.

\subsection{Location of the path and coordinate moments}

The following estimate improves the fixed annulus test to a shrinking
annulus. It concerns the ordinary observation path averaged under
$\mathbf P$; it is not an assertion about every external field.

\begin{lemma}[The observation annulus]\label{lem:observations-annulus}
There are $a,\xi>0$ and a sufficiently small fixed upper scale
$p_{\mathrm{top}}>0$ such that
\begin{equation}\label{eq:source-7}
 |\alpha_p(h_p)|\le C k_p^{-a}p^2
 \qquad(r\le p\le p_{\mathrm{top}})
\end{equation}
outside a set of $\mathbf P$-probability at most $C\exp(-n^\xi)$.
On disorder events of probability arbitrarily close to one, the same
whole-path statement holds under $Q_\mu$, with a possibly smaller $\xi$.
\end{lemma}

\begin{proof}
Under $\mathbf P$, the field in spectral coordinates has exactly its
spherical distribution. Generate the spherical spin from independent
centered Gaussian coordinates $y_i$ of precision
$D_i'=d_i+b_n n^{-2/3}$, conditioned on $\sum_i y_i^2=n$; the choice of
$b_n$ is that of Proposition~\ref{prop:edge-conditioned-law}. Before this
conditioning, write
\[
 h_{p,i}=B_{p,i}y_i+B_{p,i}^{1/2}z_i,
 \qquad m=p^2+b_n n^{-2/3},
\]
where the $z_i$ are independent standard normals. On an active coordinate
$B_{p,i}=m-D_i'$, and direct expansion gives
\[
 \mathbf E_0\bigl(\|h_p\|^2-m^2\|y\|^2\bigr)
 =-m^2\Tr R_{p,b_n n^{-2/3}}.
\]
Here $\mathbf E_0$ denotes the unconditioned Gaussian expectation. The
subtraction of $m^2\|y\|^2$ removes the large contribution of the leading
coordinates: its coefficient on the standardized $y_i^2$ term is
$(B_{p,i}^2-m^2)/D_i'=D_i'-2m=O(p^2)$.

As a quadratic form in the standardized $(y,z)$ variables, the centered
quantity has squared Frobenius norm at most $Cp^4k_p$ and operator norm at
most $Cp^2k_p^{1/3}$. Indeed the cross coefficients are
$O(p^3/\sqrt{D_i'})$, while the spectral estimates give
\[
 \min_iD_i'\ge c n^{-2/3},\qquad
 \sum_{d_i\le p^2}(D_i')^{-1}\le Cnp.
\]
The inactive coordinates are bounded by the capped second resolvent trace.
The Gaussian quadratic-form exponential-moment formula therefore bounds
deviations larger than $p^2k_p^{9/10}$ by $C\exp(-k_p^{a_0})$ for some
$a_0>0$.

This estimate survives conditioning on $\|y\|^2=n$. To see this directly,
use the same exponential tilt in the Markov bound, with its parameter at
most one eighth of the inverse operator norm. The tilted joint Gaussian
covariance, and hence its $y$-marginal covariance, remain comparable to
the original ones. A fixed number of leading covariance eigenvalues bounds
the density of $n^{-2/3}\|y\|^2$ from above. Its original density at
$n^{1/3}$ is bounded below on the chosen spectral set, by
Proposition~\ref{prop:edge-conditioned-law}. The ratio of norm densities
therefore costs only a constant.

After conditioning, the deterministic target for $\|h_p\|^2$ is
$\ell_p(b_n n^{-2/3})^2\asymp p^2k_p$. The deviation just proved is a
relative $O(k_p^{-1/10})$ error, so it first places the inferred saddle
inside any fixed surrounding annulus for all sufficiently large $n$.
On that annulus,
the derivative of $\ell_p(\alpha)^2$ is comparable to $k_p$.
Consequently the preceding deviation bound implies
$|\alpha_p-b_n n^{-2/3}|\le Cp^2k_p^{-1/10}$. Since
$b_n n^{-2/3}/p^2=O(k_p^{-2/3})$, this proves the fixed-scale assertion.

For the whole path take a sufficiently fine inverse-polynomial mesh in
$p^2$. The functions $B_p$, the saddle targets, and their relevant
derivatives have polynomial bounds; the positive-part operation defining
$B_p$ is continuous at coordinate activations. Conditional on the spin,
each coordinate of the noise has independent Gaussian increments. A maximal
Gaussian estimate in every mesh interval, followed by a union bound over
coordinates and intervals, makes their oscillation smaller than the
polynomial tolerance just used. The fixed-scale exponential estimate
absorbs this union bound because $k_r=n^{3\delta}$. Finally,
Lemma~\ref{lem:edge-biased-law}, Markov's inequality, and $L_\circ\ge1/2$
outside a vanishing disorder exception transfer the result to $Q_\mu$.
\end{proof}

The signal-plus-noise representation also gives the moment estimates
needed to test the field against individual edge coordinates. Every
$u_{p,a}$ has moments of every fixed order bounded uniformly in $a,p$ under
$\mathbf P$. For each fixed finite set of leading indices and every fixed
$j<\infty$,
\begin{equation}\label{eq:source-8}
 \|u_{p,a}-x_a\|_{L^j(\mathbf P)}
 \le \frac{C_j}{pn^{1/3}},\qquad r\le p\le p_{\mathrm{top}}.
\end{equation}
Indeed the signal is $(B_{p,a}/p^2)x_a$ and the noise standard deviation is
at most $(pn^{1/3})^{-1}$. For a leading coordinate,
$|d_a|/p^2=O((pn^{1/3})^{-2})$. The same estimates hold in $L^j(Q_\mu)$
with tight disorder-dependent constants, simultaneously in $p$ in the sense
of moment bounds at each scale: apply the signal-plus-noise formula to the
tight quenched moments of $x_a$ furnished by the biased first moment and
the concentration of $L_\circ$.

\subsection{Two replicas and the smallest scale}

In the next lemma the field is deterministic in spectral coordinates.
Consequently $E_U$ means an independent Haar average at that fixed field,
whereas $\mathbf E$ subsequently averages a random field with its
spherical observation distribution.

\begin{lemma}[Replica comparison]\label{lem:observations-replicas}
Fix a bounded range for $\|h\|/(p^2M_p)$ and take the upper scale
sufficiently small. Uniformly for fields in this range,
\begin{equation}\label{eq:source-9}
 E_U L_p(h)^2\le \exp(C+Cnp^4).
\end{equation}
For a sufficiently large fixed $C_1$, the contribution to this normalized
second moment of cube pairs satisfying
$|n^{-1}X\cdot X'|>C_1p$ is at most $C\exp(-a k_p)$.
At $p=r$, the stronger estimates
\begin{equation}\label{eq:source-10}
 \begin{split}
 E_U(L_r-1)^2&\le Cn^{-1/4},\\
 E_U\!\left[L_r^2\|m_r-m_r^{\mathrm{sp}}\|^2\right]
 &\le C M_r^2 n^{-1/4}
 \end{split}
\end{equation}
hold for the full mean vectors, written in spectral coordinates. The first
bound also holds for the uncapped zero-field likelihood $L_\circ$.
The pair-tail assertion at scale $r$ also holds in that uncapped case.
\end{lemma}

\begin{proof}
Given the overlap $\rho=n^{-1}X\cdot X'$, the Haar angular integral of a
cube pair equals that of a spherical pair. The difference is entirely in
their overlap distributions. The replica Gram comparison of
\cite[Section~15.4, Equation~(15.10)]{CM} bounds the ratio on overlap cells
of width $2/n$ by $\exp(C+Cn\rho^4)$. Above $C_1p$ and below a small fixed
cutoff, the spherical pair tail costs $\exp(-a n\rho^3)$; choosing the
cutoff small pays for the comparison on all dyadic bands. At overlaps
bounded away from zero, the normalized zero-field comparison has a strictly
negative exponential rate, by
\cite[Section~15.4, p.~118]{CM}. A bounded
field changes logarithmic weights by at most $Cnp^{5/2}$, which preserves
that rate when the upper scale is small. Below $C_1p$ the comparison factor
is $\exp(C+Cnp^4)$. This proves \eqref{eq:source-9} and its pair-tail
assertion. The uncapped argument uses the corresponding uncapped spherical
tail from \cite[Section~15.3]{CM}.

We give the finer cell estimate, now taking $p=r$ throughout the remaining
comparison. On $|\rho|\le C_1p$, Stirling's formula and the spherical
overlap density show that the binomial mass of a cell and its spherical
mass differ by a relative error
\[
 O(np^4+p+n^{-1}).
\]
It remains to compare the angular integrands within one cell. Write a
spherical pair on an orthonormal sum-and-difference frame $(q,q')$. For
$A=W_p$ their tilt exponent is
\[
 \frac n2\bigl((1+\rho)q^{\mathsf T}Aq
 +(1-\rho){q'}^{\mathsf T}Aq'\bigr)
 +\sqrt{2n(1+\rho)}\,h^{\mathsf T}q.
\]
Its oscillation on a cell is bounded by
$C(|q^{\mathsf T}Aq-{q'}^{\mathsf T}Aq'|+p^{5/2})$.
This is uniformly bounded, and its average against the normalized tilted
spherical pair law on these cells is $O(p)$. For the latter assertion,
the exact identity
\[
 q^{\mathsf T}Aq-{q'}^{\mathsf T}Aq'
 =\frac{2Y^{\mathsf T}AY'
       -\rho(Y^{\mathsf T}AY+{Y'}^{\mathsf T}AY')}
       {n(1-\rho^2)}
\]
and $\|A\|\le C$ bound the diagonal terms by $C|\rho|$. It remains to show
\[
 E_{\mathrm{sp}}|Y^{\mathsf T}AY'|\le Cnp.
\]
At a bounded cap field, apply exponential Markov with parameter $cp^2$
to either sign of $Y^{\mathsf T}AY'$. The saddle means have squared norm
$O(np)$, the covariance trace-square is $O(n/p)$, and a sufficiently small
$c$ leaves a fixed precision margin. The Gaussian logarithmic moment
generating function is therefore at most $Ck_p$. Conditioning both norms
costs a constant by the joint norm-density bound in
\cite[Equation~(15.6)]{CM}. Integrating
the resulting tail, beginning at a sufficiently large multiple of $np$,
proves the displayed estimate. For uncapped zero field, freeze a Gaussian
parameter equal to a small fixed multiple of $p^2$. The same moment
generating function bound applies, and the deficit-filling lower density
estimate in that subsection costs at most $\exp(Ck_p)$; this gives the
same integrated bound.

For completeness, the passage from this averaged oscillation to comparison
of angular integrals does not require a pointwise $O(p)$ derivative bound.
The oscillation of the exponent on a cell is bounded by a fixed constant,
so the mean value theorem bounds the difference of its exponentials by a
constant times its integrated absolute derivative. The spherical overlap
density varies by a bounded factor on a cell. Summing over cells is thus
bounded by the preceding $O(p)$ integral. Slightly enlarging the overlap
cutoff pays for boundary cells by the pair tail.

At $p=r$, $np^4+p+n^{-1}=O(n^{-1/3+4\delta})=O(n^{-1/4})$.
The normalized pair integral consequently equals $1+O(n^{-1/4})$,
which proves the first part of \eqref{eq:source-10}, since $E_U L_r=1$.
Apply the same comparison with the test $n\rho$. Its absolute value on
the retained cells is at most $Cnp=CM_p^2$, and its change within one cell
is bounded by $2$. Hence
\[
 E_U\|L_rm_r\|^2
 =\|m_r^{\mathrm{sp}}\|^2+O(M_r^2n^{-1/4}).
\]
The first-moment identity $E_U[L_rm_r]=m_r^{\mathrm{sp}}$ turns this into
an $L^2$ bound on $L_rm_r-m_r^{\mathrm{sp}}$. Finally,
\[
 L_r(m_r-m_r^{\mathrm{sp}})
 =(L_rm_r-m_r^{\mathrm{sp}})-(L_r-1)m_r^{\mathrm{sp}},
\]
and $\|m_r^{\mathrm{sp}}\|\le CM_r$ proves the second part of
\eqref{eq:source-10}.
\end{proof}

\begin{corollary}[Cube edge marginals]\label{cor:observations-static}
Jointly with the spectral variables, in the sense of
Remark~\ref{rem:edge-tight-convergence}, the quenched cube integrals of
continuous cylinder functions of polynomial growth converge to their
integrals under $\Pi$. In particular, all fixed moments of any fixed finite
set of scaled edge coordinates converge.
\end{corollary}

\begin{proof}
In the uncapped comparison, insert a bounded Lipschitz cylinder function
of the scaled spectral coordinates in each replica numerator. Changing the
overlap within a cell moves each scaled coordinate by $O(n^{-5/6})$; the
preceding weight comparison therefore remains valid. The variance over Haar
frames of the numerator normalized by $Z_{\mathrm{sp}}$ tends to zero.
Together with $L_\circ\to1$ and
Proposition~\ref{prop:edge-conditioned-law}, this proves convergence of the
quenched cube cylinder integrals to $\Pi$. Polynomially growing continuous
tests follow by the higher moment bounds and truncation.
\end{proof}

\begin{corollary}[Mean moments and the initial score]\label{cor:observations-means}
For every fixed $j<\infty$, after reducing the fixed upper scale if needed,
\begin{equation}\label{eq:source-11}
 \mathbf E\bigl(\|m_p\|/M_p\bigr)^j
 +\mathbf E\bigl(\|m_p^{\mathrm{sp}}\|/M_p\bigr)^j\le C_j.
\end{equation}
At the smallest scale,
\begin{equation}\label{eq:source-12}
 \mathbf E\frac{\|m_r-m_r^{\mathrm{sp}}\|^2}{M_r^2}
 \le Cn^{-1/20},\qquad
 \mathbf E\|\nabla_{u_r}\log L_r\|^4=o(1).
\end{equation}
\end{corollary}

\begin{proof}
Lemma~\ref{lem:observations-annulus} confines the random field to a bounded
range except with power-exponential probability. At a fixed such field,
the Haar-averaged pair-tail numerator divided by the squared spherical
partition function is at most $Ce^{-ak_p}$, by
Lemma~\ref{lem:observations-replicas}.
Under the biased frame law $L_p\,dU$, the set
$L_p<e^{-ak_p/4}$ has mass at most $e^{-ak_p/4}$. On its complement,
the expectation of the posterior pair-tail probability is at most
$e^{ak_p/4}$ times that pair-tail integral. A further Markov bound gives
an exponentially small posterior pair tail outside an exponentially small
biased exception. Since
$\|m_p\|^2=\mu_{p,h}^{\otimes2}[X\cdot X']$, this implies
$\|m_p\|^2\le Cnp$ there. The deterministic bound $\|m_p\|\le\sqrt n$
pays for all exceptions and proves the cube part of
\eqref{eq:source-11}. The spherical part follows by the same argument.

For the first estimate in \eqref{eq:source-12}, split according to
$L_r\ge n^{-0.12}$. On this event \eqref{eq:source-10} bounds the biased
second moment by $CM_r^2n^{-0.13}$. The complementary biased probability
is at most $n^{-0.12}$; Cauchy--Schwarz and the fourth moment in
\eqref{eq:source-11} give $CM_r^2n^{-0.06}$. The field exceptions are
negligible. Finally
\[
 \nabla_{u_r}\log L_r
 =n^{1/3}r^2 P_r(m_r-m_r^{\mathrm{sp}}),\qquad
 n^{1/3}r^2M_r=n^{(5/2)\delta}.
\]
Interpolation of the decaying second moment with an arbitrarily high fixed
moment from \eqref{eq:source-11} proves the fourth-moment assertion.
\end{proof}

\subsection{Likelihood dissipation and covariance errors}

The smallest-scale comparison controls a single observation. At larger
scales the useful quantity is the score accumulated along the path.
The Gaussian observation construction is a form of stochastic localization
\cite{Eldan2013}; here, conditional on the disorder, the precision follows
the deterministic spectral schedule $B_p$. We derive its likelihood identity
below and combine it with the replica estimates just proved. Write
$a_p^{\mathrm{err}}=P_p(m_p-m_p^{\mathrm{sp}})$.

\begin{lemma}[Integrated score]\label{lem:observations-score}
There is a fixed $\tau>0$ such that, on each scale interval of a fixed
bounded ratio about $q$,
\begin{equation}\label{eq:source-13}
 \int\|a_p^{\mathrm{err}}\|^2\,d(p^2)
 \le Ck_q(q^\tau+k_q^{-\tau})
\end{equation}
outside a set of $\mathbf P$-probability at most $C\exp(-n^\xi)$.
The bound holds simultaneously on a geometric list of intervals, with
fixed enlargements, and also holds in expectation.
\end{lemma}

\begin{proof}
Use $t=p^2$ as the channel parameter and take the filtration generated by
the eigenframe and observations through time $t$. In particular, every
disorder event is measurable at time zero. Between coordinate activations,
$\dot B_t=P_p$, and the Gaussian likelihood identity gives
\[
 d\log L_p=a_p^{\mathrm{err}}\cdot d\mathcal B_t
       +\tfrac12\|a_p^{\mathrm{err}}\|^2dt,
\]
where $\mathcal B$ is the innovation Brownian motion in the active
subspace. The bracket of the martingale is the score integral. There is
no singular term at an activation because the added precision starts at
zero. These statements also follow by It\^o differentiation of the
posterior partition function, including the compensating interaction
change $-dB_t$.

Put $H_q=Ck_q(q^\tau+k_q^{-\tau})$, with $\tau$ small and fixed.
At the first endpoint, the likelihood-density identity gives
$\mathbf P(\log L_p<-H_q)\le e^{-H_q}$. At the last endpoint,
\eqref{eq:source-9} gives
\[
 \mathbf P(\log L_p>H_q,\ h_p\text{ in the bounded range})
 \le \exp(-H_q+C+Ck_q q).
\]
Choose $C$ large and use Lemma~\ref{lem:observations-annulus} outside
the range. An exponential martingale inequality now bounds the event
that the bracket exceeds a sufficiently large multiple of $H_q$ while
these endpoint inequalities hold. This proves
\eqref{eq:source-13}. The logarithmic number of intervals is absorbed
by the power-exponential bound. Both means have norm at most $\sqrt n$,
so the integral has a polynomial deterministic bound; integrating the
exception proves the expectation assertion.
\end{proof}

Let $\mathcal G_n$ be the intersection of the disorder events on which
the fixed-tolerance annular covariance bounds of
Proposition~\ref{prop:edge-cap-inputs}, the static value and projected-mean
comparisons of Lemma~\ref{lem:edge-static-interpolation}, and the associated
projection tails hold on the full scale and fixed field ranges used below.
All tolerances and ranges are fixed before taking $n\to\infty$; their
finite intersection has probability tending to one.
The next estimate is asserted on $\mathcal G_n$; its exceptional probability
is measured before any conditioning on that event.

\begin{lemma}[Integrated covariance error]\label{lem:observations-curvature}
Put $D_p^K=p^2K_p-\Id$ on $\mathcal H_p$. After decreasing $\tau>0$ if
necessary,
\begin{equation}\label{eq:source-14}
 \int_q^{2q}\frac{\|D_p^K\|_{\mathrm F}^2}{k_p}\,d\log p
 \le C(q^\tau+k_q^{-\tau})
\end{equation}
outside an event whose intersection with $\mathcal G_n$ has
$\mathbf P$-probability at most $C\exp(-n^\xi)$. Fixed enlargements of
the interval are allowed, and the assertions hold simultaneously on a
geometric list.
\end{lemma}

\begin{proof}
On the shrinking annulus of Lemma~\ref{lem:observations-annulus}, the
spherical covariance has norm $O(p^{-2})$. By rotation invariance within
the flat cap its restriction to $e^\perp$, $e=h_p/\|h_p\|$, is scalar.
The frozen Gaussian linear-tilt bound gives an $O(p^{-1})$ error between
the spherical radial mean and the saddle mean. Dividing by
$\|h_p\|\asymp p^2M_p$ shows that the transverse covariance differs
from $(p^2+\alpha_p)^{-1}\Id$ by $O(p^{-2}k_p^{-1/2})$.
The remaining radial covariance is rank one. Thus its contribution,
and the shrinking saddle parameter, are already bounded by the right
side of \eqref{eq:source-14}.

It remains to compare $K_p$ and $K_p^{\mathrm{sp}}$. The spherical
posterior, evaluated along the cube observation path, has innovation
drift $K_p^{\mathrm{sp}}a_p^{\mathrm{err}}dt$. Subtracting the two
posterior mean equations gives
\[
 da_p^{\mathrm{err}}
 =(K_p-K_p^{\mathrm{sp}})\,d\mathcal B_t
       -K_p^{\mathrm{sp}}a_p^{\mathrm{err}}dt
\]
between activations. Multiply $p^2\|a_p^{\mathrm{err}}\|^2$ by a smooth
nonnegative cutoff supported on an enlarged scale interval and equal to
one on $[q,2q]$, with bounded derivative in $\log p$. It\^o's formula
and integration by parts give
\[
 \int_q^{2q}\|p^2(K_p-K_p^{\mathrm{sp}})\|_{\mathrm F}^2\,d\log p
 \le C\int_{\mathrm{enlargement}}\|a_p^{\mathrm{err}}\|^2\,d(p^2)
       +C|\mathcal M|.
\]
Activations add squared new mean coordinates with nonnegative coefficients,
so they may be discarded on this side of the inequality. The cutoff is
zero at both endpoints; the displayed energy identity is applied over the
entire interval, not at a stopping time. Its martingale is
\[
 \mathcal M=2\int\chi(t)t
   \langle a_{\sqrt t}^{\mathrm{err}},
      (K_{\sqrt t}-K_{\sqrt t}^{\mathrm{sp}})\,d\mathcal B_t\rangle.
\]
For concentration, stop this stochastic integral alone at the first
failure of the annular covariance bounds, and set its integrand to zero
on $\mathcal G_n^c$, a time-zero event. Its bracket is then bounded by
$C\int\|a_p^{\mathrm{err}}\|^2d(p^2)$. The stopped integral agrees with
$\mathcal M$ on $\mathcal G_n$ and the whole-path annulus event.
The martingale inequality, Lemma~\ref{lem:observations-score}, and
Lemma~\ref{lem:observations-annulus} therefore prove the assertion.
Division by $k_q\asymp k_p$ gives the stated normalization.
\end{proof}

We finish with an angular estimate that will identify derivatives of the
limiting measure. This estimate uses the remaining Haar randomness without
restricting first to $\mathcal G_n$.

\begin{lemma}[Angular score]\label{lem:observations-angular}
For a cap field $h$, write $\widetilde h=U^{\mathsf T}h$ and
$\widetilde\phi_p(\widetilde h)=\phi_p(U\widetilde h)$.
Let $A$ be a fixed skew-symmetric matrix supported on finitely many leading
spectral coordinates, and set
$\mathcal R=(A\widetilde h)\cdot\nabla_{\widetilde h}$. There is
$\tau'>0$ such that
\begin{equation}\label{eq:source-15}
 \int_q^{2q}
 \frac{\mathbf E|\mathcal R\widetilde\phi_p(\widetilde h_p)|^2}
      {n^{2/3}p^2}\,d\log p
 \le C_A(q^{\tau'}+k_q^{-\tau'}).
\end{equation}
\end{lemma}

\begin{proof}
The coordinates on which $A$ acts belong to the cap for all sufficiently
large $n$.
By \eqref{eq:source-4}, the joint density of the frame and spectral
observation relative to independent Haar and the spherical observation
law is $L_p$. Condition on the physical capped matrix $W_p$, the physical
vector $h_p$, and the spectral coordinates of that vector. The spherical
observation density is now fixed, and $L_p$ depends only on the fixed
physical data and spectrum. The remaining rotation of the spectral frame
in $e^\perp\cap\mathcal H_p$ is therefore exactly Haar. In physical
coordinates, the differentiation direction is $UA\widetilde h_p$, which
is orthogonal to $h_p$ and has length $|A\widetilde h_p|$. The spherical
mean in the cap is radial. Haar averaging consequently gives
\[
 \frac{\mathbf E|\mathcal R\widetilde\phi_p(\widetilde h_p)|^2}
      {n^{2/3}p^2}
 \le C\mathbf E\left[
 |Au_p|^2\frac{p^2\|a_p^{\mathrm{err}}\|^2}{k_p}\right].
\]
Apply \eqref{eq:source-13} in expectation and interpolate with the higher
moments in \eqref{eq:source-11} and the coordinate moments following
Lemma~\ref{lem:observations-annulus}. H\"older's inequality on probability
times logarithmic scale gives \eqref{eq:source-15}, with a smaller positive
exponent. No event depending on the unused spectral refinement was inserted
before the Haar average.
\end{proof}
